\section{渡口、侨州与账簿没有写下的人}
\label{sec:eastern-jin-depth-migration-fiscal}

长江并不是一条把北方灾难自然隔在岸外的线。三一三年前后，江北渡口要让携带家属的官员、散兵、商人和地方武装依次过江，也要让逆流而上的盐、米、布帛和军书进入荆扬之间。船只的载重、滩口的水势和守渡者愿意放行谁，决定了迁徙不是一条整齐的南下箭头。有人带着旧籍、私家文书和可以转托的亲属名号抵达建业；也有人在淮北停下，依附坞壁或军府，等候北方局势再变。东晋最初面对的并非一批已经安置好的“侨民”，而是许多尚未决定向哪一种保护关系缴纳劳力与粮食的人。

\begin{event}{\eventanchor{JH-0313-ZU-TI-NORTHERN-EXPEDITION}{祖逖北渡与豫州屯田}}{建兴元年（三一三年）}{东晋前身与河南地方集团 · 江北渡口、豫州屯田区}{北渡、屯田与地方结盟可核；收复范围、兵粮数和归附程度不能精确复原}
祖逖率部北渡时，建业一方尚未成为正式帝都。渡江之后，最要紧的并不是在地图上画出一条前线，而是让跟随者有地可耕、有粮可食，并同河南已有的居民和武装谈出可以共处的次序。屯田在史书中常以成效概括，落到当地却包含开垦、守护水源、分配耕作收益和防范邻近势力的连续安排。祖逖的军队能够向北行动，正因为它不完全等待一条自南方中枢送来的粮道；这也决定了它对将领个人的关系网高度依赖。

这种依赖给建业带来缓冲，却不是可复制的北伐制度。祖逖所能联系的坞壁、流民和旧吏未必接受同一套命令，军屯也会把耕作者、军人和地方豪强卷进新的负担分配。祖逖死后，原来由个人威望和局部供给维系的网络难以无损交接。北渡的价值因此不只在收复了多少城镇，而在于它暴露了一项长期矛盾：朝廷若要重新进入河南，必须提供比名义任命更稳定的粮饷和保护；若做不到，前线便会成为地方集团自行谈判的空间。
\end{event}

侨置州郡正是在这种不稳定中成为一套可用、却并不完整的行政语言。把原籍的州郡名称带到江南，可以保存仕宦资格、编入某种名册，也能让朝廷把北来官员安排进官号与俸禄的秩序。但地名的迁移不等于土地、人口和税源同时迁来。一座侨郡的长官可能领有旧地的名义，而他实际面对的是江南本地村落、早到的移民、军府和地主的交错网络。对于被征发的人，决定生活的往往不是郡名来自何处，而是谁来催租、谁来征兵、发生纠纷时能向哪一处县署求助。

江东旧有的水田、湖泊、圩垸和市镇，也不因北方人抵达而成为无主的容器。朝廷需要北来官员的文书技能和军事经验，本地精英则掌握乡里关系、田地信息和水路通行的知识。两者的合作可以进入婚姻、任官和军府，也会在田宅、佃客和治安上产生摩擦。后世门第叙事容易让人只看见几个高门的相互接纳；从地方尺度看，更多人经历的是租佃关系如何改变、劳役由谁摊派、陌生的军队能否驻进村落。传世材料很少直接留下这些人的连续声音，不能把沉默补写成一致的欢迎或抵抗。

建康的财政也不是一座都城打开府库就能解决的问题。朝廷可以授予官职、发出诏令，却要经过州郡、军府和地方中介才能得到谷物、布帛、船只与人手。江淮之间一有战争，粮船便要避开劫掠与封锁；荆州军镇的态度又影响上游物资能否顺江而下。因此，东晋的“南方”从来不是一个均质后方。建康、京口、武昌、江陵和会稽所处的水陆节点不同，向中枢提供的资源和提出的条件也不同。王敦、苏峻等人能以外镇压力进入都城，正因为他们并非凭空拥有军队，而是掌握了让人、船和粮沿江移动的一部分接口。

将迁徙只写成士族南渡，会把这些接口的代价遮住。官员的迁居需要仆从与依附者，军队的移动需要挑夫、舟子与沿线供给，侨置的文书需要地方吏员逐项执行。每一次看似抽象的“安置”，都可能改变一户人家是否离开故居、是否把子弟送入军府、是否承担额外运输。墓志与家传主要记录有能力留下姓名的家庭；它们能让我们看见籍贯如何被保存，却不能替全体迁徙者说话。正因如此，东晋早年的连续性应理解为许多局部关系被勉强接续，而不是一个北方国家整体搬迁到江南。

北方的变化持续进入这套南方账目。洛阳、关中、河北的政权换手会改变流民从哪里来，也会改变建康是否必须增兵淮河。朝廷为了留住将领而给出的名位与土地安排，可能减轻眼前威胁，却同时提高地方军府的独立性。侨州郡解决的是身份与任官的一部分难题，不能替代田亩、赋役和军事保护的实际分配。东晋最早的一百年，便是在这一连串不完全的接合中维持下来：渡口仍然开放，文书仍能传递，某些军队仍有粮食；但每一项都依赖具体地方的合作。

\begin{sourcenote}[南渡的名目不能换算成总人数]
《晋书·元帝纪》《王导传》《祖逖传》与《资治通鉴》可以校核出镇、北渡和朝廷建立的顺序；侨州郡记录、墓志与城址则提示迁徙留下的局部痕迹。它们不构成一次人口普查。流亡、死散和户口数字多带叙事或行政目的，本节只据此讨论迁徙怎样进入渡口、任官、屯田与赋役，而不将其换算成精确的南渡规模。
\end{sourcenote}

\section{襄国、邺城与凉州：北方的征发怎样落在城外}
\label{sec:eastern-jin-depth-northern-regimes}

黄河以北的政治重组并不只发生在帝号更替时。襄国、邺城、长安、姑臧和龙城都是人们需要缴纳粮食、寻求保护、传递消息的地方；它们之间的力量差异首先表现为谁能守住城门、修复道路、把征来的物资送到军队手中。后赵的扩张使河北出现一个强有力的中心，但强中心也意味着更多地方被纳入军役、运输和宫廷营造的要求。史书对石勒、石虎的记载充满人物褒贬，不能把这一叙述直接化为居民的共同感受；仍可看出，河北城镇与乡村被卷入了比建康朝廷更直接的北方军政竞争。

\begin{event}{\eventanchor{JH-0328-SHI-LE-CAPTURES-LIU-YAO}{洛阳附近俘获刘曜}}{太和三年（三二八年）}{后赵与前赵 · 河南、洛阳周边道路}{刘曜被俘及前赵主力崩解可靠；战场位置、兵数和败因有争议}
石勒与刘曜争夺的不是一块抽象的中原。洛阳附近的道路连接河内、关中和河北，城镇、渡口与仓储决定一支军队能否把胜利变成继续行动的能力。刘曜被俘后，前赵的政治中心失去统率者，许多原先依其保护而行动的人必须重新判断向谁交粮、向谁请援。后赵由此获得河南的战略入口，但接收一地并不只是夺取城池：旧吏、守军、商旅和农村中介都要在新的命令链中寻找位置。

此役的后果是权力天平向襄国倾斜，而非北方秩序立刻稳定。前赵旧部的去向、关中与河南每座城的服从程度，无法由胜负一项推定。对住在道路旁的人来说，较快到来的变化可能是军队索取粮草、逃亡者增多和市场通行受阻。正是在这些不易进入帝纪的日常压力中，新的政权才逐渐显出轮廓。
\end{event}

石勒称帝后，襄国需要把军事胜利转化为常态化的命令。官号、礼制和都城营建能吸引文武集团，也能告诉地方中介该向何处呈报；它们却不自动消除从军者、牧地控制者、旧郡县吏和新附豪强之间的利益差异。后赵建立较完整的官僚和军政体系，是理解其力量的必要起点；把制度名称直接等同为全境均匀执行，则会忽略每一座城、每一条粮路的差别。河北的地力和人口使襄国能调动大规模资源，也使宫廷的征发更容易转化为地方负担。

石虎掌权后的继承与扩张尤其显示这种代价。宫廷内部的废立固然决定谁能发布命令，命令能否抵达远方又取决于地方军队、城防和运输。史书常用残暴或奢侈概括石虎时期，读者不应把道德标签当作解释。更具体的问题是：为了维持都城、战争和宗室竞争，哪些州郡被要求提供劳力，哪些运输线被军队优先占用，地方中介如何在服从与自保之间选择。现有材料无法逐县复原答案，却足以提醒我们，后赵的繁盛与紧张同出于它对河北资源的集中调用。

后赵崩解时，邺城的争夺把先前被压住的矛盾放大。石虎死后，继承竞争使将领、宗室和不同地方网络不再听命于一个中心；冉闵在邺城的行动又引发大规模暴力。史书中关于“羯”“胡”“汉”的分类、杀伤数字和群体行为，均有强烈的政治叙事与后出整理，不能把它们转写为固定族群之间的天然战争。可以确认的是，政权崩解使许多居民和军人失去原有保护，城内外的暴力与迁徙骤增；不能确认的是每一群人的人数、动机和受害范围。把这一危机放回邺城、道路和征发网络中，才不会把人群名称误作完整解释。

凉州的经验又有所不同。河西走廊依靠绿洲、灌溉、城郭和西向道路维系生活，前凉能够在中原反复动荡时保持相对独立，不是因为它远离政治，而是因为当地的官吏、豪强、商旅和军队能围绕这些节点形成自己的资源安排。前秦于三七六年灭前凉，改变了最高政权的归属，却未使绿洲立即成为关中行政的复制品。谁来接收仓储、怎样使用旧官、寺院与商路能否继续供养旅行者，都是征服之后才开始的事务。

河西同时提醒我们，十六国不能只从洛阳或建康看。使节、僧侣、商人和军队通过走廊来往，带来不同语言、经卷、技术和政治消息。文书、题记与遗址能照见其中一些时间点，却保存得极不均衡；一份契约或一个洞窟题名不能代表整条走廊。与正史并读时，它们让人看见，所谓边疆并非中原命令抵达的终点，而是多种区域网络交错、也可能拒绝单一中心的地方。

慕容氏在辽西和东北的发展同样建立在区域道路上。棘城与龙城连接辽西、幽州和高句丽方向，前燕能够利用城镇与地方关系抵抗后赵压力，却也必须不断处理部众、迁入者和邻方之间的变化。远征丸都等事件的胜败可大体辨认，行军路线、俘获人数和居民反应却不应被叙事夸张填满。东北的政权并不是中原强弱的附录；它改变了河北军队的后方与选择，也让当地社会在多个索取资源的中心之间反复调整。

\section{成都的水路、关中的粮道与北伐的边界}
\label{sec:eastern-jin-depth-campaigns-finance}

东晋的北方目标常以“恢复旧土”概括，但真正限制行动的是水路、山道和可被持续供给的军队。桓温由荆州向西进入巴蜀时，利用的是长江上游与支流构成的通道；进入成都之后，面对的却是如何让新任官员、旧有吏员和地方豪强在同一套税役安排中继续工作。成汉的灭亡让建康重新取得益州的名义与资源，却并未让成都平原自动变成可随时调用的北伐仓库。上游的粮食、船只和人手首先要在地方社会中被征集，再通过多个军政层级转为远征能力。

\begin{event}{\eventanchor{JH-0354-HUAN-WEN-NORTHERN-EXPEDITION}{桓温入关中而后撤军}}{永和十年（三五四年）}{东晋与前秦 · 荆州、汉水、关中外围}{北伐、进至关中与因粮运撤退可靠；兵数、战况和责任分配有争议}
桓温北伐前，收复成都使荆州军府拥有更大的声望和上游纵深。军队沿汉水及关中道路推进时，越靠近前秦核心，越不能依赖后方的临时接济。山谷、渡口和城市仓储既可能帮助行军，也可能在敌对力量与季节变化下成为阻断。前线即使取得局部胜利，粮食仍要从远处运来；一条被拉长的补给线会迫使将领在继续进攻、驻守城邑和保存军队之间作出取舍。

撤军并不是“意志不坚”的同义语。前秦尚未稳固，给东晋留下了机会，也使关中每一处接应都充满不确定。桓温不能把一支依靠荆州与江南供给的军队无限期留在陌生地区，尤其当当地的合作、收成和道路安全都不能保证时。北伐因而成为一种政治资源：它能在建康提升统帅声望，向士人展示恢复旧都的承诺；同时又暴露中枢不得不依赖大军府的财政和人力。军功越大，朝廷对掌兵者的制衡反而越困难。
\end{event}

三五六年桓温进入洛阳、修复晋帝陵，提供了北伐最具象征性的时刻。帝陵的修复使旧朝记忆与新的军功在一处地点相遇，却不能代替稳定的治理。洛阳周围的村落、城防和运输尚在反复易手，东晋无法仅凭一场仪式长期驻守。三六九年枋头失利则把问题写得更清楚：当军队远离水陆供给节点、需面对北方政权的集结时，河道、粮运与后方协调会比“收复”的口号更早决定行动边界。战报中的兵数与奇谋颇多异文，不宜用来制造精确战场；前后两次行动共同说明，象征性进入与持续控制是两件不同的事。

前秦在关中逐步集中资源，也是在这种限制下显得强大。苻健、苻生到苻坚的权力更替，王猛的整顿，以及对前燕、前凉、代国的连续征服，使长安能比东晋更直接调动北方的粮食、兵员与城镇。政策能实行到何种深度必须逐地检验，不能以“集权”二字概括；但前秦确实把多处道路与资源节点纳入较统一的军事计划。这种集中既提供南伐的条件，也把不同地区的兵源和供给压在一套可能过度延伸的命令系统之下。

梁、益北部州郡的反复易手，尤其能看见地方社会如何承受大政权的边缘竞争。对建康而言，它们是汉水、蜀道和北伐之间的接口；对住在当地的人而言，战争意味着守城、避兵、运输和重新确认向谁纳税。史书写“取”“失”往往只用一两个动词，实际接收却包括旧军队如何处置、地方官是否留下、田地和仓储由谁管理。将这些地区只看成东晋与前秦的边界，会遮住它们自身在多次征发下形成的脆弱性。

军事财政也塑造了建康内部的权力。一个能够筹集船队、安置军士和控制江面的人，能在朝廷有危机时迅速变成不可忽略的政治行动者。桓温的声望使中枢既需要其防卫北方，又担心其挟军功干预继承与任官。东晋并非没有文官体系，问题在于文官诏令必须经过地方供给网络才能变成行动。北伐之所以一再发生又一再难以延续，正因为它同时满足了皇统的象征需要、军府的地位需要和流亡士人的期待，却很少拥有足以跨越淮河、汉水与关中长期维持的财政底盘。

\begin{causalchain}[北伐为何反复越过边界又难以留下]
收复巴蜀为荆州军府提供上游纵深与政治声望；关中远征却把军队拉入更长的运输线。短期内，北伐能吸纳流亡者的期待并压制敌方边境；成本则由沿江、沿汉的运输者、田地与地方中介承担。撤军留下的不是单纯的失败，而是军府权威继续增长、中枢更依赖掌兵者、下一次北伐更难摆脱同一供给难题的反馈。
\end{causalchain}

\section{长安的译场、河西的洞窟与地方文字的尺度}
\label{sec:eastern-jin-depth-religion-regions}

四世纪末到五世纪初，战争并没有使道路只供军队使用。僧侣、经卷、施主、抄写者与商旅也在关中、河西、江南和西南之间移动，只是他们留下的材料与帝纪不同。佛教的传播不能被理解为一个皇帝下令、全境便共同信仰的过程；一座译场、一则造像题记或一处石窟遗址，分别照见不同人群在不同时间作出的供养与书写选择。它们把政治中心之外的知识、手工与交通带入叙事，同时也要求读者承认可见者常是有资源留下文字的人。

\begin{event}{\eventref{JH-0401-KUMARAJIVA-CHANGAN}{鸠摩罗什抵达长安与后秦译经}}{弘始三年（四〇一年）}{后秦 · 长安与关中道路}{抵长安和译经活动可靠；入城日期、译场规模及姚兴参与程度有争议}
姚兴支持鸠摩罗什在长安译经，使一座由后秦控制的城市成为跨区域知识网络的节点。译场需要的不只是君主的态度：来自西域的僧侣、懂得校读与汉文书写的人、纸墨、居处和日常供养，都要在关中城市中得到安排。经卷因此并非从宫廷单向流出，而是在僧侣、弟子、施主和抄写者的关系中逐步形成可以携带的文本。长安的政治保护降低了这一网络的风险，却不能说明城市居民都参与译场，更不能把一部经论的完成日期等同于它在各地被阅读的日期。

译经活动也改变人们看待城市资源的方式。一个能够供养僧侣和抄写的中心，需要较稳定的粮食、纸墨与住宿安排；战乱时期，这种稳定本身就是一种稀缺的政治能力。后秦以此吸引文教资源，和它以军队守卫关中并非两条互不相干的线。可是赞助并不等于控制思想，僧侣之间、不同译本之间以及后来读者之间都可能有不同理解。把姚兴的支持写成国家替所有人决定信仰，既超过材料，也抹去译场中实际协作的人。
\end{event}

法显自长安西行、经河西与西域而归的经历，则让人看见另一种道路使用方式。他的文本记录的是求法者的选择与见闻，不是沿途所有居民的社会调查；其中有关国家、风俗和距离的叙述也须结合成书语境阅读。但它提示河西的城镇、寺院与交通并非只在中原征服时才出现。旅行者能否继续前行，取决于绿洲补给、地方保护和同行网络，任何一项断裂都可能改变路线。宗教移动由此把地方政治、商路与文本生产连在一起，而不能被缩成“文化交流”的抽象箭头。

炳灵寺的题记与河西石窟材料提供了不同于正史的时间尺度。题名、年号和供养者姓名能帮助辨认某一地点的营造与政治语汇，却不能让我们从一处洞窟推出整片地区的信仰分布。洞窟开凿可能早于或晚于一块题记，供养者的身份也未必代表附近农户、军人或商人的共同选择。它们的珍贵之处，是让宗教活动落在岩壁、河谷、工匠和资金的具体条件上：石窟不是凭空出现的，需要组织劳力、获得材料并长期维护。

宁州的《爨宝子碑》又将视角转向西南地方。碑文保存官职、家族与书写样式，使一个远离建康与长安的地区在可核日期中露出轮廓。它首先是家族的公开陈述，不是全体地方社会的名册；对官职、世系和政权关系的说法都需同其他材料比对。即便如此，碑石仍提醒我们，三、四世纪的区域生活不只由都城与战役构成。地方精英通过文字处理身份，地方道路把云南、巴蜀与更广的网络相连，政权更替未必以同一速度进入每一个山谷和聚落。

江南的宗教实践也在迁徙与军府压力中改变。寺院能够提供结社、供养与文本流通的空间，士人和地方富户的参与会让它们进入城市和乡里的资源网络；但现存僧传与碑刻偏向著名人物和成功机构，不能把它们转化为普通家庭的统一信仰史。道教网络在会稽、浙东等地同地方社会和政治动荡的关系同样复杂。孙恩起事后来被官方语言归入特定宗教或人群类别，这有助于理解朝廷如何命名威胁，却不能替代对饥馑、征发、海路与地方武装的分析。

宗教材料的限制反而使它们更有用。正史让人追踪一座都城何时易主，题记、墓志、经录和旅行文字则提示同一年中不同地方仍在修墓、刻石、抄经、行旅。它们不能填满普通人的沉默，却阻止我们把所有人都写成随朝廷命令移动的棋子。长安的译场、河西的洞窟、宁州的碑刻和江南的寺院，共同显示多政权时代的文化不是统一帝国遗留的余波，而是在道路、供养和地方选择中不断重新结成的网络。

\section{淝水之后：胜利不能接收北方，军功也不能安置江南}
\label{sec:eastern-jin-depth-fei-liuyu}

前秦南伐之前，苻坚已经把关中、河北、河西和代北等不同区域纳入扩张网络。它的优势在于能从更大的范围调兵、筹粮和发布命令；弱点也在于这些部众、道路和地方中介并不因一次征服就获得同样的忠诚与补给能力。东晋的北府兵则依托京口、淮河和江南较近的水陆节点，规模与构成不能由后世名称想象成固定军团，却能在本土附近较快取得协同。淝水的关键，不是神奇口令使庞大军队瞬间消失，而是两种资源组织方式在一条前线上的碰撞。

\begin{event}{\eventref{JH-0383-FEI-RIVER}{淝水之战与前秦南伐崩解}}{太元八年（三八三年）}{东晋与前秦 · 淮河、淝水及北返道路}{东晋取胜与前秦军崩溃可靠；兵数、口令和伤亡数字有争议}
淝水附近的水道要求军队在渡河、列阵、传令和补给之间保持秩序。前秦远道而来，前线每一次移动都会牵动后方粮运与不同部众的协调；东晋军则更靠近能够补充人马和船只的江淮节点。战斗的具体瞬间被后世故事反复塑造，不能逐帧还原。可以把握的是，一旦前秦阵势与指挥失去稳定，远征军便很难在陌生且拉长的道路上重新聚集；苻坚北返，原先被征服或暂时归附的地方力量也开始重新计算自己的选择。

东晋保住江南，并没有获得一把接收北方的钥匙。河北的慕容垂、北地的姚苌以及各地旧部在不同条件下建立新中心，前秦的瓦解更像多处道路与城镇失去共同协调，而非一片土地等待南方军队进入。建康没有足够稳定的粮饷和军镇体系去同时接收这些地区；它能庆祝胜利，却仍要面对淮河防务、荆州军府和地方财赋的长期压力。淝水的政治意义正在于阻止一次南方覆亡，而不是完成北方统一。
\end{event}

北魏兴起提供了另一个后果。拓跋珪在代北重新聚合力量、迁入平城，说明北方的整合将围绕新的城郭与征发网络展开。平城的仓储、道路和官署使军事集团能够同农耕城镇、工役和新附人口建立更持续的关系；这不等于旧部落关系立刻消失，也不等于制度在所有地区均匀落实。它只是给北方提供了一个能把人马和物资重新集中起来的中心。淝水后真正出现的不是东晋北上，而是多个区域在不同时间、以不同资源重新组织。

东晋内部的危机也没有因胜利终止。会稽与沿海的孙恩起事、桓玄从荆州东下，均使江道、海路、城门和外镇成为决定朝廷存亡的接口。官方叙述会把起事者归为叛乱或宗教性群体，文本中的称呼不能直接替代其成员的动机。战争、征发、土地关系和地方保护的破裂，都可能促使不同人进入同一支武装。建康反复需要外部军队救援，意味着皇统仍有吸引力，却也意味着掌兵者可以借救援重新分配中枢权力。

刘裕在这一环境中成长。平定桓玄、控制京口与江上通道后，他北伐南燕、入广固，随后又处理卢循与谯蜀等危机。每一次军事成功都依赖不同的水陆条件：山东方向要跨淮河并围攻城郭，岭南与江南的危机牵涉海路和长江，巴蜀接收要面对三峡与成都平原。把刘裕的行动概括为个人雄略，会遮住其军政网络如何吸收地方供给，也会遮住被征发者承担的运输、守城与安置成本。

四一六年至四一七年，刘裕入洛阳、长安，后秦覆亡，一度使东晋军重新站到关中核心。可是关中并不会因入城而自动提供稳定后方。刘裕南返后，赫连勃勃的夏夺取长安，显示当地军队、道路与政治关系尚未被东晋持续接合。前线的失去不能只归咎于一位统帅离开：江南、河南与关中之间的供给跨度，地方接应的有限性以及军镇本身的利益，都使占领难以变成持久治理。北伐再次带来巨大的声望，却没有改变东晋财政和区域控制的根本限制。

四二〇年受禅时，建康以晋朝的礼仪语言把掌军者转为新皇帝。这个仪式能解决名义上的继承，却不能让侨州、军府、江东旧族和北来家族立刻服从同一套日常安排。对许多地方而言，改变首先表现为新的官号、征发命令和保护关系，而不是都城仪式本身。东晋的终结因此并非故事的终点：它把南方军政网络交给刘宋，同时把未能解决的地方财政、军镇与北方多政权竞争推入南北朝的新阶段。

\begin{debate}[淝水胜利与“统一机会”]
《晋书》及《资治通鉴》保存东晋获胜、苻坚北返和前秦分裂的年序，却不足以证明东晋拥有立即统一北方的物资条件。研究可以讨论前秦崩解是否创造机会，但不能把机会误作已经具备的接收能力。东晋军队的供给、地方合作与外镇政治，正是判断这一问题不可跳过的条件。
\end{debate}

\section{邺城失序以后：城门、田地与不能化约的人群}
\label{sec:eastern-jin-depth-ye}

后赵在石虎死后迅速陷入继承战争，显示一套依靠宫廷威望和大规模征发维持的统治，未必能留下可以平稳转交的次序。邺城的宫门、军营和仓库是争夺的焦点，但城外的乡村、来往于河北的道路和守卫各处城镇的部队同样决定了谁能活用这一中心。继承者需要军队承认，军队需要粮食，供应粮食的地方又要判断新命令是否能保护自己。朝廷内的废立因此很快变成更大范围的迁徙、掠夺和重新结盟。

冉闵在邺城的崛起常被放在一个极端简洁的民族叙事中。传世史书对于群体名称与杀伤数量的记述，带着胜者、后继政权和道德判断层层加工；所谓“羯”“胡”“汉”并不是能用来精确划定每个人立场的名册。仍不能回避的是，后赵崩溃期间发生了严重的针对性暴力，许多人因军队、城门与道路失去保护而被迫逃亡。材料不能让我们替每位受害者确定身份和去向，恰好说明不要以一个整齐的群体故事覆盖不同城镇中不相同的灾难。

邺城的混乱并没有只留下废墟。对于周围的农户、坞壁与小城而言，新的问题是播种能否继续、谁来征收粮食、逃入城内是否比留在村里安全。战争会毁坏收成，也会让拥有武装与储粮的人获得谈判优势；地方精英可能组织自保，普通家庭则可能在逃亡、依附或交出劳力之间做选择。史书极少让这些选择以第一人称出现，但从政权反复争夺城镇、征发与安置的记载可以知道，北方社会并不是在一场宫廷政变后静止等待新王朝。

前燕在这一缝隙中向中原推进。慕容儁将政治中心由辽西带入河北，利用的既有慕容集团的军事关系，也有后赵崩解后城镇对稳定保护的需求。把前燕的扩张写成一个固定“族群”压倒另一个群体，会忽略它必须接收旧吏、安排守军、维持田赋与道路的现实。前燕取得邺城等地，意味着河北的资源重新进入一个政权的征发范围；同时也意味着新中心必须面对前政权留下的军人、流民和地方中介，不能只靠国号完成接收。

前秦在关中的崛起与前燕并行。苻健建立前秦时，长安周边的农业、关隘和西向交通给了政权较厚的基础；苻坚即位后，王猛等人推动整顿，意在削弱强宗与调整官吏。传统叙述很容易把这种过程写成一位能臣立即令天下整齐，实际执行仍须经过郡县、仓储和地方合作。诏令可以指明国家希望怎样征税、任官和惩治，不能单独证明每一处都按同样方式转变。前秦的力量正在于能不断把区域资源重新接入长安，而不是已经消除了区域差异。

三七〇年前秦灭前燕，是这一能力的高点，也把新的困难带入关中政权。前燕原有的官员、军队和河北城市不可能因战败自然改用长安的节奏。接收者需要决定谁留下、谁迁徙、粮仓归谁、边境由谁防守；被接收者则要判断何种服从可以保全家业与人身。史书中的降服名单和胜利数字容易压缩这种过程。较可把握的是，前秦由此得到更广的兵源与道路，同时把异质地区纳入一个需要持续协调的国家。后来南伐的规模，正来自这一集中；其脆弱性也埋在同一集中之中。

东北与河北之间的战争还影响了更远处的家庭。辽西、幽州和高句丽方向的道路不仅运送军队，也传递逃亡者、消息和货物。前燕对丸都的进攻、同周边势力的往来，在不同史书中有不同叙述重心；我们不能由一份胜利叙事推断所有迁徙或地方服从。可以确认的是，这些地区并不被动接受中原命名。城郭、山地通道、河流与地方首领的关系，使它们在前赵、后赵、前燕与后来的北魏之间具有自己的行动空间。

当北方在邺城、长安和辽西不断重组时，建康看到的往往是北伐的机会与威胁。可是南方朝廷很难把北方某一政权的崩溃直接换成自己的控制范围。要进入淮北或河南，仍须筹粮、安顿前线、取得地方接应。北方的失序确实让逃亡者和消息进入江南，也使建康的士人更执着于恢复旧土；但它同样提醒人们，每一次大政权覆亡之后，出现的首先是多种地方选择，而不是等待某一“正统”接收的空地。

\section{前秦的道路网络：从河西到淮南的过度伸展}
\label{sec:eastern-jin-depth-former-qin}

前秦的扩张不是一连串孤立的胜仗。它把关中、河北、河西、四川北缘与代北方向的城镇、仓储和军事人群逐步接入长安的命令。三七三年取梁、益北部州郡，三七六年灭前凉、并吞代国，都是将道路与资源节点纳入更大范围的过程。对长安而言，这些地区可以提供军队、牲畜、粮食和战略纵深；对当地人而言，最高政权变化首先表现为官员更替、运输要求、守军驻扎和旧有保护关系的重新估价。

河西被纳入前秦时尤其值得放慢叙述。姑臧及周边绿洲并非一块单纯的“边地资源”，其灌溉、城市生活、寺院、商路与西域联系构成了长期形成的区域秩序。征服军进入后，要面对的不只是前凉君主是否降服，还包括谁管理水源、谁保管仓粮、谁维护向西的通行。前秦的纪年材料能确定大致的政治转折，却不能让我们断言每一座绿洲在同一年、以同样方式完成行政改造。后来吕光在凉州建立新的政权，也说明远方地区能够利用城镇和道路重新组织，而非只被长安单向支配。

前秦对前燕的征服同样带来接收问题。河北人口与城镇可以扩大南伐的动员能力，却把不同军政传统、地方中介和刚被征服的群体带入同一支远征军。苻坚能够命令各地集结，不表示行军途中每一队人马都有相同补给、语言或对战争目标的理解。大规模军队越过淮河前，粮食要从多个区域汇集，文书要穿过漫长道路，地方长官还要在本地生产与国家征调之间分配稀缺物资。规模是优势，也使任何一处误判更难被及时修正。

东晋面对前秦压力时，不能只靠都城礼仪回应。京口的北府兵、淮河沿线的守备、荆州的兵力与江南的粮运需要同时被接合。谢玄在京口训练军队，意味着朝廷试图把北来军人、地方供给与边防节点组织成较可调用的力量。北府兵不应被浪漫化为一个天然同质的集团：其成员、兵额和供给在材料中并不完整。它的重要性在于显示东晋的防卫能力来自一个具体区域的军政经营，而不是凭“南方富庶”四字自动产生。

前秦南伐前的外交与心理压力也有物质的一面。边境城镇要判断是否坚守，沿路州郡要准备粮运，江淮居民则面对可能到来的征发与战乱。史书里的诏书与言辞保存统治者希望别人如何理解局势，却不能直接说明每一个地方怎样行动。有人可能因前秦强盛而观望，有人因旧有关系支持东晋，还有人首先考虑如何避开军队。正是这种多层选择，使战争不能被归结为两个国家在地图边界上的单一碰撞。

淝水前秦军的崩溃因此并不神秘。远距离调兵的系统在平时能够显示威力，在前线一旦出现命令混乱、部众离散或补给受阻，修复的成本便格外高。东晋军的胜利使这个系统失去可信度，许多先前被连接到长安的地方网络便有理由重新寻找自己的中心。慕容垂、姚苌等人的行动不需要被解释为一个预先写定的“民族复兴”；它们是在前秦权威下降后，利用各自拥有的部众、城镇与道路所作的具体选择。

淝水后的前秦遗产并未消失。长安、邺城、姑臧、平城等地仍保有行政经验、人口与交通价值，后来的政权会以不同方式接收它们。北魏的平城经营、后秦对长安的利用、河西各国的延续，都说明一场大战无法抹去城市和地方社会积累的能力。战争改变了谁能征发，却没有让灌溉停止成为灌溉、让渡口停止成为渡口。多政权时代的连续性，常常正藏在这些比国号更耐久的设施与劳动中。

\section{会稽、广州与荆州：江南地方生活并不都在建康城内}
\label{sec:eastern-jin-depth-southern-local}

东晋后期的江南不能只由建康和北伐构成。会稽的海岸与水网、广州通向海路的港口、荆州控制长江中游的军府，都有不同的财富、交通和社会关系。中央需要它们提供税赋、船只和兵员；地方家庭则在田地、盐业、渔业、手工业、市场与军役之间安排生活。传世史书对都城政变和著名将领保存较多，对普通人怎样承担赋役保存较少，不能因此把地方社会写成没有历史的背景。

孙恩起事从会稽及沿海扩展时，官方文本会以“妖贼”“道民”一类名称归类行动者。这些词能够告诉我们朝廷如何识别威胁，却不能替代对地方压力的解释。沿海和岛屿提供航行与暂避的空间，水网也使军队难以像在平原上那样迅速封锁。参与者中可能有因赋役、战乱和保护关系受挫而行动的人，也可能有利用宗教网络、家族关系和海路的地方武装；现有材料不足以把他们归为一个动机一致的群体。重要的是，东晋的国家能力在此受到检验：它必须调动地方船只、军队和粮食，才能重新连接那些看似属于它、实际却有自身网络的区域。

卢循由岭南及海路北上，又把广州的区域性凸显出来。岭南同建康相距遥远，海路既是贸易、迁徙和信息通道，也能在危机中成为军队移动的走廊。朝廷对广州的任命不等于立即控制每个港口与乡里；当地方军事力量取得船只、补给和同伴时，海上距离反而可以削弱中枢反应。刘裕等人平定这一危机，依赖江防和水军的组织，却不应被写作一次纯粹的都城胜利。沿海居民所经历的征船、避兵、贸易中断和战后重新纳入地方秩序，才是事件延续到日常生活的方式。

荆州的意义又不同。它处在长江中游，能够控制来自巴蜀的物资与通往北方的水陆路线，因而军府长官常有介入建康政局的实力。桓玄从荆州东下并夺取朝廷，说明都城若无法独立掌握沿江供给，皇统会成为外镇争夺的资源。桓玄的覆亡也并非恢复一条自然秩序：救援者、地方将领与军府都会因勤王获得新的位置，下一轮权力集中由此开始。刘裕的崛起正建立在这种军事救援与地方网络重新排序之上。

地方生活并不只在危机中可见。墓志、碑刻和城址偶尔留下家庭怎样书写籍贯、婚姻、官职与安葬；这些材料多来自有资源的群体，不能代表佃客、奴婢或军户。它们仍能纠正一种误解：江南社会不是北方士人抵达后才开始存在。北来者要在已有村落、市场、水利与宗族关系中安顿，地方人也会借新朝廷的官职、婚姻和军府重新安排自己的上升渠道。所谓南北融合不是一个瞬间，而是一连串以土地、身份与劳役为媒介的协商。

宗教网络在这一协商中有时提供可见的组织方式。寺院、道教社群和地方祭祀可能成为供养、结社、文本流通或互助的节点，也可能被官府在危机中视为需要控制的联系。不能把一则僧传或叛乱指控变成群众信仰统计，但可以看到，宗教语言与实际的道路、船只、家族和粮食相互交织。它既不是脱离政治的精神背景，也不等同于政治动员的单一工具。

东晋末年的地方危机说明，江南的稳定从来不是自然产物。朝廷能在建康举行朝会，仍须依赖会稽、京口、荆州、广州及无数州县把谷物、布帛、船只与人力接入网络。一处外镇失控、一条水路被截断或一次地方征发失衡，都可能让都城的命令停在纸上。刘裕后来能够取代晋室，正因为他逐步掌握了这些接口中的一部分；而新王朝仍不得不面对同样的问题：如何让军功所建立的控制转为不完全依赖战争的日常治理。

\section{从广固到长安：刘裕的远征与接收的限度}
\label{sec:eastern-jin-depth-liuyu-campaigns}

刘裕北伐南燕时，山东并不是一片可由建康随意进入的旧土。军队需要跨越淮河，沿途寻找补给，并面对广固这样有城防与储备的中心。南燕控制山东的一部分道路与居民，东晋军的到来会迫使当地人重新判断守城、提供劳力或避入他处的风险。围攻广固的结果可由传世材料确认，兵员数量、每一段路线和城内情形则不宜写成全知的战场。较有解释力的是，刘裕能把水陆调动、军府组织和北伐声望结合起来，使一次远征成为建康政治中的筹码。

南燕覆亡后，东晋并未因此掌握整个北方。山东的接收需要安排守军、官吏和粮运，北魏及其他北方势力仍在改变更广范围的力量关系。刘裕返回南方处理卢循、谯蜀等问题，也表明统帅的资源始终有限：同一支军政网络不能无限分布在山东、岭南、巴蜀和建康之间。战争的成效会把更多地区拉进征发体系，也会增加维护这些地区的成本。对地方社会而言，新的旗帜到来通常带来的是新的驻军和文书，而非立刻确定的和平。

四一六年北伐后秦，东晋军通过河南进入关中。洛阳与长安具有巨大的旧都象征，但远征的困难仍是具体的。黄河、道路、城镇仓储和当地合作决定军队能否继续推进；由江南向关中延长的补给线，则要求后方保持相对稳定。后秦覆亡让刘裕达到前所未有的军功高度，却没有消除长安与建康之间的距离。当地的守备与政治关系尚待重新组织，任何被忽略的节点都可能成为新的力量进入的入口。

赫连勃勃建立的夏趁东晋军主力南返夺取长安，最能说明接收的困难。东晋获得关中不是没有意义，它改变了当地行动者的选择，也让刘裕的声望进一步上升；但占领没有形成能独立维持的地方网络。将长安失去只写成刘裕个人离去，反而把问题说小了。关中需要持续的粮饷、能够同地方合作的官员、足以应对邻近政权的守军；这些条件并不能由一场胜利或一份任命自动产生。

刘裕回到建康后，军功、朝廷人事和对外镇的控制使受禅成为可能。晋恭帝退位并非社会生活突然改写的时刻。侨州郡仍在、江东旧族仍在、北来家庭和军府仍在，地方税役和水路运输仍要靠既有中介处理。受禅仪式的作用，是让新的掌权者使用能够被文官、军人和地方精英理解的正统语言；它解决了“谁能代表朝廷”的问题，不能立即解决“谁能让各处按日供给、守法和服役”的问题。

这也解释了为什么东晋到刘宋的转换应当被看作接收，而不是断裂或自然延续。新王朝继承了建康的官僚、江南的财赋和刘裕的军事网络，也继承了北方难以久守、外镇需要约束、地方社会对征发敏感的旧问题。北方多政权的持续存在使南方始终不能把安全只寄托于长江天险；而江南的水网、城市与人力则使南朝能够在一次次危机后重新建立中心。连续与改变，恰在这些具体的道路、田地、船只和名册里同时发生。

\section{国号之间的日常财政：谁替军队找到粮食}
\label{sec:eastern-jin-depth-fiscal-local}

在东晋与十六国的叙事中，财政很少以一部连续账簿的面貌保存下来。正史记得诏令、赏赐、赈济和大规模征发，却很少逐村列出田亩、租额与运输损耗。材料的缺口不应使财政退成背景。每一个能够保持军队的政权，都必须把粮食、布帛、牲畜、铁器与劳力从地方带到城郭和前线；每一次战争，又会把更多家庭卷入耕作中断、运粮、修城和避难。我们看不见完整数字，仍能从道路、仓储、屯田、侨置和军府的运行看见资源怎样成为政治。

建康的收入依赖江南州郡，却不能把“江南富庶”当成一个永远有效的答案。稻作、水利和水网确实提供了较强的生产与运输条件，然而水路也意味着粮船容易受季节、战事与地方封锁影响。中央发出的征调要通过刺史、郡守、县吏、豪强与军府才能落到土地和船只上。中间每一层都有自己的信息和利益：地方官要维持治安，军府要先保军粮，乡里要避免把存粮全部交出。朝廷能否得到足额供给，取决于这些关系是否仍能合作，而不只取决于诏书写得多么严整。

侨置给财政带来的问题尤为复杂。迁来的官员和家族需要籍贯、居住与仕宦资格的安排，原有居民则担心新增的军役与赋役如何分摊。侨州郡的名目能帮助国家识别某些人，却不能让田地、人口和税源自动对应。若一个地方同时有实土县与侨县，谁负责征收、谁承担徭役、遇到诉讼向何处申报，都可能成为长期的实际难题。传世材料不允许我们对每处安排下结论，但它们足以否定“南渡者一到便自然占有江南资源”的简化图像。

北方政权的财政同样不是单一模式。后赵以河北为重心，能借密集城镇与农业区征发人力，也要负担不断的军队调动、宫廷竞争和城防需要。前秦以关中为基础向外扩张，吸收前燕、前凉等地区后获得更多资源，同时必须让远方道路不断把物资输送到可能发生战争的边缘。河西的绿洲、辽西的牧地与城镇、关中的农业区各有不同供给方式；将其统称为“国家财力”，会遮住各地实际承受的征发和协商。

军屯提供一种在战时把粮食与军队绑在一起的办法，却不是没有成本的技术。祖逖经营豫州时，屯田使前线不至完全依赖江南远运；参与耕作与守卫的人必须在生产和战斗之间分配时间，周边居民也会面对土地、水源和保护关系的重新划分。后来各政权在边地与军镇的安置，常有类似目标：让军队就地获得部分供给，减少长线运输。可是屯田能否成功，要看水利、治安、地方接应与征发强度，不能由制度名称推断。

货币和市场也在战争中受到影响。史书中关于赏赐、赎买或物资价格的零星记载，难以拼成完整的经济统计，却提醒我们人们并非只以谷物生活。城镇需要盐、铁、布帛、纸墨与牲畜，军队需要马匹、车船和修具；商路的畅通会影响一座城能否渡过围困。政权更替可能改变关卡和征收者，但很少能让交换需求消失。河西与江南的商旅、寺院与市场，都是将区域生产接入更大网络的地点，也可能在军队经过时首先承受搜括与中断。

对普通家庭而言，财政不是抽象的“国用”，而是何时被要求交粮、谁来验收布帛、家中劳力是否被带去运兵、战乱后能否回到原有田地。此类经验极少留下直接记录，尤其是没有墓志与文书的农户、军户和依附者。我们不能替他们编出统一的情绪，却应把这种不可见性保留在叙事中。大政权的扩张若不接到劳动与运输的代价，便会显得像地图颜色自然伸展；真正的国家能力恰恰由无数这样难以在正史中署名的人支撑。

宗教机构也进入资源流动。译场需要纸墨与供养，寺院需要土地、施主、工匠和维护，造像或开凿洞窟需要长期组织劳力。王室或地方豪强的赞助能为这些活动提供条件，却不应被直接解释为全民信仰或完整的经济控制。题记所见的供养者姓名，多让我们看见有余力留名的人；没有留下名字的劳动者和周边住户仍在材料之外。把宗教放回财政与地方生活，并不是把信仰简化为利益，而是承认文本和建筑都需要具体的物质条件。

多政权时代最显著的财政事实，或许正是资源的重复索取。一个地区在政权易手时，旧税役未必立即消失，新守军又可能要求粮草与劳力；逃亡者进入新的县境，也会改变当地的负担与保护关系。最高权力的更替对地方而言不是一个日期，而是一段重新确认谁有权征发、谁能提供安全的过程。东晋与十六国的长期并立因此并非只造成政治分裂，也让地方社会不断在不同的账目、名册和军队之间寻找生存余地。

\section{长江以北的边缘不是空白：淮河、汉中与山地城镇}
\label{sec:eastern-jin-depth-frontiers}

淮河与汉水之间的地区在东晋叙事中常被称为北伐前线，仿佛它们只在南方军队越境时才有历史。事实上，渡口、堤防、城镇、市场和山路早已让居民、地方武装与官吏形成自己的生活节奏。西晋崩溃后，北方诸政权、东晋军府和地方集团都试图把这里变成缓冲或通道；住在其中的人则必须在不同保护者之间衡量风险。边界不是一条静止的线，而是一系列谁能通过、谁能储粮、谁能守住夜间城门的具体问题。

淮河水系使军队的速度和补给高度依赖季节。枯水时可渡之处增多，雨季与洪水又会改变船只和道路的选择；一座城若控制渡口，便能影响人、粮和消息向南北流动。东晋在京口和江淮布置军队，不只是为了象征收复，也是为了防止敌方在接近江南的节点站稳。前秦南伐时，淮河与淝水的空间条件把大军拉进狭窄的行动范围；淝水后，东晋虽守住南岸，却仍须长期经营这些水陆接口。

汉中与梁、益北缘把关中、巴蜀和荆州连接起来。山地道路使一支军队难以像在平原那样迅速转向，也使熟悉地形的地方力量具有重要性。桓温进攻成汉、北伐关中，以及前秦取得梁益北部，均显示这些地区不能被简单称作“通往某地的走廊”。城镇、关隘、粮仓和居民的合作决定走廊是否真正开放。一次军事进入若不能留下守备和供给，便只是暂时穿过，而不是治理。

巴蜀的山地与成都平原之间也存在这种不对称。成都拥有农业和城市资源，外围道路却可能在战时迅速成为限制。成汉的政权依靠区域资源建立，继承冲突又削弱其协调能力；桓温入成都后，东晋得到上游支点，却必须重新处理益州的任官、税役与旧军队。对于当地家庭而言，谁在建康称帝或襄国称王未必是最直接的变化，驻军是否到来、山道是否安全、市场能否恢复才会先进入日常。

代北与河西的边缘性又与江淮不同。代北的牧地、山谷与营地为拓跋集团提供行动空间，平城的建立则把这些军事网络接入城郭、仓储与官署。河西依靠绿洲间的距离、水源和西向商路，能够在中原政权更替时保持区域联系。后来的北魏、北凉和其他政权都要面对同一个事实：边缘不是等待中心填补的空处，而是拥有资源、交通与地方中介的复杂区域。中心若不能理解这一点，扩张便会在道路的尽头变成过度伸展。

“边民”这样的称呼也必须小心。官方文书和史书常以此分类需要征发、安置或防范的人群，却不能据此知道每个人的语言、身份或政治期待。某些人可能同时与城镇市场、牧地网络和军府发生关系；在政权变化时，他们的选择也未必与官方标签一致。把族名或边地称谓当作稳定本质，会重演史书为治理而作的分类，反而失去理解地方社会如何行动的机会。

边境宗教与文化的材料进一步显示这种交错。河西洞窟、墓葬、文书和僧侣行旅留下的痕迹，往往同时指向地方供养、西域联系和政权年号。它们不能为我们画出精确疆界，却能说明一条道路在政治地图上属于谁，与人们实际怎样往来并非同一问题。东晋和十六国时期的区域史必须同时看见两者：政权为争夺城镇与关隘而战争，普通的通行、交易、求法和迁徙又可能越过这些临时边界继续发生。

\begin{textwindow}[道路材料不是精确行军图]
《水经注》及后出的地理叙述能帮助辨认河流、城址与道路关系，但成书年代、所引旧闻和文本层次并不相同。它们适合提示淮河、汉水、河西与代北的空间限制，不足以逐日复原一支军队的路线，更不能把后出的地名关系直接投回四世纪的每一年。
\end{textwindow}

\section{材料如何改变读法：从帝纪到石刻的交叉}
\label{sec:eastern-jin-depth-evidence}

东晋与十六国留下的材料极不均匀。帝纪和载记使人较容易追踪都城易手、称号更改与著名战争，却多由较晚的编纂者以晋朝或后继政权的视角整理；十六国原史散佚，后代辑佚保留片段，常无法恢复原有上下文。若只依这些文本，读者很容易以为历史只由皇帝、将领和几个都城推动。墓志、碑刻、造像记、寺院遗址、城址和旅行文本提供不同尺度的校正，却各有选择性，不能拼成一部无缝的地方档案。

王兴之夫妇墓志的价值正在这种有限性上。它以明确的姓名、日期、官职与安葬关系，让建康附近一个有资源留下石刻的家庭出现在我们面前。它不能告诉我们全城迁入多少北方人，也不能让我们推知没有墓志的军户、佃客与本地家庭如何生活。但若没有它，朝廷传记中的“侨姓”便更容易成为抽象类别。墓志让读者看见迁徙、婚配、仕宦和葬地怎样在一个家庭的自我呈现中相连，也提醒我们每一份可见证据都有其不可见的周边。

《爨宝子碑》与炳灵寺题记也应这样使用。碑刻能校正地方官职、家族和书写的时间，题记能把供养活动与某种年号、地点联系起来；二者都不能代表宁州或河西的全部人口。洞窟的开凿、题名的刻写、政权的改元可能并不同时，不能把一个日期扩大为整座寺院或整片区域的政治归属。材料越具体，外推的边界反而越需要清楚。

佛教文本同样有自己的时间。鸠摩罗什抵长安、译场活动和经卷流传可以分别在不同材料中得到支持，但一部文本的翻译、抄写、校订和后世流布不是一个时刻。僧传会突出高僧的德行与行迹，经录关心文本传承，政权叙事则可能强调王室赞助；它们回答的问题不同。把它们并读，可以看见长安如何成为知识与供养的节点；若将它们混成一条“全民信佛”的证据链，反会失去材料各自的力量。

对于战争，材料边界尤其影响解释。淝水、参合陂、广固或长安的胜负可以大体确认，兵数、伤亡、临阵口令和每支部队的路线却经常服务于军功、败责或后世文学想象。谨慎不是让战争变得没有意义，而是把意义放回可以核验的条件：远征军的补给、城郭的守备、河流与道路、地方合作以及政权权威的变化。少用一个不可靠的精确数字，反而更能看见为什么一场胜利会改变后续选择。

族群称谓也需要同样的处理。正史中的“胡”“羯”“氐”“羌”“鲜卑”等词可能指向政治身份、军事集团、地域来源或作者的分类习惯，并不等于现代意义上边界清晰的民族共同体。某一政权使用某种名号，也不能说明所有被归入其中的人共享同一行动理由。把称谓放回城镇、军队、婚姻、征发和外交的语境，才能避免让后出的分类替代当时人复杂而变化的关系。

来源意识不是把每段正文变成出处清单。它应当改变读者读事的方式：当材料只允许确认称帝或入城时，便不虚构宫廷心理；当一方传记给出惊人的兵数时，便不把它当作统计；当一块碑石照见一户人家时，便不让它代表整座城市。东晋与十六国的历史之所以值得这样阅读，正因为其多政权、迁徙与地方生活不能由一部连续官方档案完全覆盖。承认证据的空白，不是削弱叙事，而是让道路、家庭、城市与战争在它们真正能被看见的尺度上出现。

\section{从国号到地方：十六国并非一张整齐清单}
\label{sec:eastern-jin-depth-polity-reading}

“十六国”是后世为方便叙述而形成的总称，它不能替代当时人的政治地理。不同政权出现、改名、迁都和灭亡的时间并不整齐，另有许多短暂或区域性的行动中心难以被这个名称完全容纳。前赵、后赵、成汉、前凉、前燕、前秦、后秦、后燕、南燕、北燕、北凉以及代北的北魏前身，各自在不同的城镇、道路和资源条件中形成。将它们排成一列国号，能够帮助读者辨认年序；若把国号当作边界明确、内部一致的国家，便会错过多政权时代最重要的现实：一座城、一个军府或一条商路常常比抽象疆界更能决定人们的选择。

成汉的历史提示区域政权不能被当作中原崩溃的附属现象。李氏在成都建立中心，依托的是巴蜀的农业、山地通道与地方军政关系；继承冲突削弱了这种协调，桓温的军队才得以沿水陆路线进入成都。成汉覆亡后，巴蜀并未从历史中消失，而是重新进入东晋的任命、税役与上游防务。政权变更改变了谁在成都发布命令，却没有解除山路、河流、地方豪强和旧吏对新政府的限制。

前凉与后来的河西政权则提示，地方延续可以比中原朝廷更长。姑臧和各处绿洲通过灌溉、市场、寺院与西向交通维持区域生活，既能接受来自关中的官号，也能在远方中心衰弱时重组自己的军事与行政安排。前秦的征服、吕光的据凉、北魏的接收，都不是简单的“中央到边疆”过程。每一次接收都要面对水源、仓储、旧官、商旅和宗教机构，任何一项都可能让纸面上的隶属关系在当地呈现不同面貌。

前燕、后燕与北燕使东北和河北的联动进入年序。慕容集团能从辽西进入河北，靠的不只是军事行动，还有城镇、道路与地方中介在后赵崩解后的重新选择。后燕在淝水后兴起，又在北魏压力下失去优势；北燕在龙城重组，说明东北并非在一个大政权倒下后就自动归属下一个强者。高句丽方向、辽西道路与地方城郭给这些行动者提供了不同于关中和江南的资源条件，因而也产生不同的政治时间。

后秦的长安与北魏的平城进一步说明，都城不是地理上的终点，而是一项持续建设。姚兴能在长安支持译场，是因为后秦保持了关中城市的供养与保护能力；拓跋珪迁入平城，则把代北军事网络同仓储、官署和城郭相连。两座城市都吸引文书、僧侣、工役与军人，却各自面对不同的地理与人群关系。把它们都归入“北方政权”，并不能说明它们怎样取得粮食、怎样安置新附者，或为何其文化与军政安排会不同。

东晋同这些政权的关系也不能缩成“南北对立”。建康同前赵、后赵、前秦、后秦、燕系政权之间有战争、使节、流民、边境贸易和关于旧晋名义的复杂互动。南方朝廷会把北方恢复视作政治目标，北方行动者却未必以同样的正统尺度组织自己。对于河南、淮北、关中和山东的地方社会，重要的不只是建康是否宣称收复，更是眼前哪支军队能守城、哪一方征粮较重、哪条道路仍能通行。多政权的意义正在于，政治合法性、军事控制和日常生活的保护关系经常并不重合。

因此，阅读这一时期不宜问“哪一个国家代表中国”便停止。更有用的问题是：某一政权在何处能持续征发，哪些人和城镇帮助它传递命令，它的扩张在何处因补给或地方合作而停顿，政权崩解后又是谁接住道路、军队和人口。这样读，前赵的洛阳、后赵的襄国与邺城、成汉的成都、前秦和后秦的长安、燕系的龙城、北魏的平城，以及东晋的建康，便不再只是更替的都名，而成为不同社会被组织、被征发并重新寻找保护的具体地点。

\section{章节收束：南渡不是退场，北方也没有一次“失去”}
\label{sec:eastern-jin-depth-conclusion}

从永嘉危机到刘裕受禅，东晋与十六国的年序始终由两种运动交错推动。一种是政权、军队与家族沿道路、水道和城镇不断移动：洛阳、长安、襄国、邺城、成都、姑臧、龙城、平城与建康依次或同时成为新的中心。另一种是较慢的重组：田地重新耕种，侨置重新编入名册，军府寻找粮食，寺院与译场获得供养，地方家庭在战乱后选择安葬、婚配、迁居或依附。前者容易进入纪年，后者常只在墓志、题记、文书和沉默中露出边角；二者却不能分开。

南渡使晋朝名义在建康继续存在，但它不是北方旧社会的完整搬迁。江南原有的水利、市场、村落与精英关系给了新朝廷资源，也要求北来者与之协商。侨州郡、军府和任官制度提供了接合工具，却不能抹平赋役、土地和治安上的差异。每当建康需要北伐或应付外镇，地方资源便被重新拉进政治中心；每当征发过重或保护失效，会稽、荆州、岭南与江淮又会显示它们不只是被动后方。

北方也不能被理解为东晋从地图上“失去”的一大片土地。刘渊、石勒、慕容氏、苻氏、姚氏、拓跋氏及许多地方行动者都在不同资源条件下建立中心。战争带来可怕的暴力和大规模流动，但城镇、绿洲、农田、关隘和道路并未消失；它们成为下一轮政权能够出现的条件。淝水击败前秦，阻止了东晋的覆亡，却没有把北方交给建康。北魏在代北和平城的经营、后秦在长安的文化赞助、河西与东北的持续重组，正是这一未被接收的空间如何继续生成历史的证据。

东晋末年的刘裕将长期积累的军功转成受禅，显示南方朝廷已经能从军事网络中产生新的皇统。然而他的关中远征也表明，进入旧都与长期经营旧都并不相同。刘宋继承的不是一个已经统一的南方或一个等待收复的北方，而是一套仍要处理军镇、地方财政、迁徙家庭和边境道路的关系网。由此开始的南北朝，并非突然换页，而是东晋与十六国时期各项压力在新的政权名称下继续展开。

读者若在这一时期感到国号繁多，并非因为历史缺少主线。主线正在于：每一次国号的出现，都要回答谁能保护城镇、谁能调动粮食、谁能让文书与人穿过道路；每一次国号的消失，又会把这些问题交给新的行动者。连续的年序使我们看见这一接力，地方材料和来源限度则提醒我们，不要把无法被完整保存的人群、家庭与劳动，从这条主线中删除。

\section{看见城与看不见的人：地方生活的最后一层}
\label{sec:eastern-jin-depth-local-lives}

政治年序常以都城陷落、皇帝即位和将军出师为节点，地方生活的时间却未必同步。洛阳失守后，远在江南的村落不会在同一天改变租佃关系；建康称帝后，淮北的家庭也不会立刻知道新的诏令是否能带来保护。人们通过亲属、商旅、军队、逃亡者和地方官听到消息，再根据水路是否安全、田地能否耕作、孩子是否会被征役而调整选择。这个过程很少被连续记录，不能用想象填补，却必须作为理解政权韧性的前提。

城镇是最容易留下痕迹的地方。城门、仓库、官署、寺院与市场把不同身份的人聚在一起，也让征发和保护更容易被看见。一座城被攻取时，史书往往记载主将、降者与日期；对居民而言，变化还包括商市是否开放、粮价如何波动、守军向谁征粮、家人是否躲入城内。城墙可以提供屏障，也会在围困中把有限粮食、疾病和恐惧集中起来。我们不能由一段攻城记叙复原全城经验，但也不应让“攻取”二字抹去这些具体后果。

乡村与坞壁的材料更少。地方武装可能保护邻里免受流兵侵扰，也可能凭借武力要求粮食与劳役；豪强、旧吏、军府和流民集团的界线会随时改变。一些人会依附拥有土地和武装的家庭，一些人会逃向山谷、江岛或另一座城，另一些人仍留在原地试图守住播种和收获。这样的选择不是宏大正统叙事的附属，而是国家能否获得税粮、军队能否获得补给的基础。正因为选择分散且材料不全，我们只能描述其条件，不能替所有地方人宣布一种共同态度。

女性、儿童、奴婢、佃客和普通军户在传世文本中尤其难以看见。墓志会记载婚姻、家世和少数女性的德行，却多为有能力建墓立碑的家庭制作；法律与诏令提到户口和役使，却反映国家希望如何分类，而不是每户怎样承受。承认这种偏差并不是放弃社会史，而是避免以少量精英材料虚构完整的日常。迁徙中的照料、家庭分离、依附关系的重组和劳役的性别分配，都是可能塑造生活的重要问题，只是现存材料难以给出统一答案。

地方文化也以不均匀的方式延续。家族可以借墓志、碑刻和谱系声明籍贯与官职，寺院借题记和经卷留下供养者姓名，市场与道路则很少保存普通交往的文字。建康的士人讨论玄学、佛学或政治，河西的僧侣与商旅传递经卷，宁州的家族刻石纪念官员；这些都是真实的文化活动，却不能被简单排列为从都城向地方单向扩散。不同地区选择、改写或忽略中央文化资源的方式，正是多政权时代文化活力的一部分。

环境与交通对地方生活的约束同样具体。江南水网带来灌溉、渔业和便利运输，也会在洪水、战时封锁或征船时成为风险；河西绿洲依赖水源，任何政治动荡都可能影响灌溉与通行；北方平原的农耕与牧地关系则要面对军队移动和城镇征发。战争并没有取消这些条件，只是让它们更容易成为争夺对象。谁控制渡口、堤防、山道或水源，谁便能影响军队与居民的选择。

将这些尺度与东晋、十六国的年序并读，可以避免两种相反的误解。一种把地方人写成只会被政权推动的沉默背景；另一种则因为材料稀少而替他们创造没有证据的完整故事。较稳妥的做法，是让城址、墓志、题记、文书和正史各自说明它们能照见的部分，同时把看不见的范围标出。如此，建康、襄国、长安、成都与姑臧的兴亡才会重新落到实际生活：有人在其中获得官职和保护，也有人失去田地、亲属或通行的道路。

这层地方生活还解释了为什么政权更替总比史书中的日期更长。新朝廷要重新登记、征收、任命和安置，旧关系不会立即停止；地方家庭要观察新守军是否可靠，新的法律是否真能到达县里。四二〇年的受禅标志东晋皇统结束，却不能决定江南每一处寺院、田庄和渡口在次日如何运作。多政权时代的历史并非由这些缓慢过程拖累，而正是由它们获得持续性。若没有耕作、运输、抄写、婚配和安葬，任何国号都无法在地图上停留太久。

\begin{sourcenote}[地方社会的可见与不可见]
墓志、碑刻、造像题记、城址与零星文书能为家庭、宗教和区域活动提供局部锚点；《晋书》《华阳国志》及编年材料则较擅长保存朝廷与战争年序。两类材料都不能替代完整的户口、赋役或日常生活统计。本节将它们并置，是为说明地方生活如何与政权竞争相连，而非把少量精英证据扩写为所有居民的共同经验。
\end{sourcenote}

\subsection{一部没有完成的接收史}

把东晋与十六国放在同一卷中，最容易看见的并不是一条由混乱走向统一的直线，而是反复发生的接收。石勒接收前赵的河南，桓温接收成汉的成都，前秦接收前燕与前凉，后秦利用前秦留下的长安，北魏再把代北军事网络接到平城；刘裕进入洛阳、长安又不得不面对无法久守的关中。每一次接收都包含三层不同速度的变化。最高层是国号、皇帝与都城的改变，史书最容易记下；第二层是官员、军队、仓储和道路的重新分配，需要数年甚至更久；第三层是家庭重新安排田地、依附、婚配和迁徙，常常只留下零碎痕迹。

这三层不能相互替代。一个政权可以在一天之内攻入城门，却不能在一天之内让旧吏相信新命令、让军粮稳定抵达、让居民忘记前一轮征发。反过来，一处地方即使更换了多次旗帜，水源、市场、坟墓、寺院与亲属网络仍会使生活保有某些连续性。连续性不等于人们免于战争；它只是说明人们必须在战争中继续耕作、运输、抚养与纪念，国家才有可能在下一年重新征发和统治。

因此，东晋的存在不能只由建康未失来衡量，十六国的力量也不能只由疆域最大时来衡量。更值得问的是，它们能否把不同区域接到一套足以运行的关系中，又在何处被粮道、军镇、地方合作和继承危机拆开。这个问题将东晋南渡、后赵兴亡、前秦扩张、淝水崩解、北魏兴起和刘裕受禅连成同一条年序：不是国号不断替换的目录，而是人、物资和制度在多种空间中被迫重新安放的历史。

这种视角也改变了“成功”与“失败”的尺度。祖逖在河南的经营、桓温的关中远征、谢玄的淮河防卫、刘裕的长安入城，都不能只用最后留下多少疆土衡量。它们分别把屯田、荆州军府、京口兵源和江南水陆供给带到极限，也使后来行动者知道哪些道路可以暂用、哪些地方必须先取得合作。前秦的失败并不否定它曾将关中、河北与河西组织成能发动远征的网络；东晋的胜利也不表示它已经拥有接收北方的能力。把这些结论放回具体地点，读者便能同时看见政权选择的创造力与它们无法摆脱的限制。

材料的沉默同样属于这部历史。没有一份完整账簿告诉我们所有粮食来自哪里，没有一份名单记录每个逃亡者的终点，也没有一种族名能准确概括在战争中行动的人。可靠的叙事不以虚构填满它们，而是让可核的都城、渡口、城墙、经卷、碑石和道路同未能保存的生活并置。这样，所谓分裂时代才不只是中心失控的阴影，也是一段区域社会在压力下不断创造新联系、又不断承受新代价的编年。

从这个意义上说，读者不必在国号繁多中寻找一条虚假的直线。顺着渡口、城镇、粮道、寺院与家庭的选择前行，便能看见同一时代的多种速度：皇统可以骤变，地方秩序只能缓慢修补；军队可以越过河流，供给和信任却未必随之抵达。
这一差异有时决定结局。

\section{江淮的防线不是一条线：船、堤与临时的安全}
\label{sec:eastern-jin-depth-river-defence}

东晋能够以长江为屏障，并不意味着江面本身提供了无需维护的安全。江防依赖船只、渡口、沿岸城镇、瞭望和能够在不同水位行动的人。建康的官员可以把长江说成天险，真正让它发挥作用的却是京口、历阳、采石等节点的守备，以及舟师、船户、沿江州县能否及时供给人力与粮食。水道在和平时连接市场与迁徙，在战争时也让敌军、流民和求援消息快速移动；它既保护江南，也把江南暴露在必须不断应对的交通压力之下。

淮河比长江更直接地显示边界的可变。河道宽窄、支流、季节和堤防会改变军队可以渡过何处、粮船必须绕开何处。一座小城控制渡口，便可以影响来自北方的逃亡者、南方的援军和地方商旅的选择。所谓守淮并不是把军队排在一条恒定岸线上，而是在多个能够通行的地点维持观察、仓储和响应。东晋的北府兵之所以重要，正因为它把军人、京口的地方资源和江淮前线接到较近的指挥网络中；这份优势也有边界，不能自动延伸到黄河或关中。

防线的日常代价落在沿岸社会。修堤、造船、运输、守城和供给驻军，都需要劳力与材料。对一户以耕作或水上生计为主的家庭，军队到来既可能带来保护，也可能意味着船只被征用、收成被优先征发，或市场因戒严而受阻。史书比较容易记录将领在何处胜败，较少记录船户是否得以补偿、农田在军营附近如何收割。不能用无根据的细节替代这些沉默，但应当认识到，边防能力是由这样的协商和损耗构成的。

孙恩、卢循等危机之所以能越出都城的预期，也与水路有关。会稽和沿海的支流、港汊、岛屿使行动者可以避开陆路大军的追击；岭南至江南的海路又使遥远地区在特定条件下与长江流域直接相连。朝廷把这些行动称为叛乱，反映了中枢的安全视角；从区域角度看，海路同时也是贸易、迁徙、宗教和地方关系的通道。军队若只按陆上郡县的逻辑理解它，便难以迅速封锁人员与物资的流动。

东晋将领在江上取得声望，往往也因此获得政治分量。能调集舟师、维持渡口、在危机中将军队送到建康的人，不只是执行朝廷命令，也掌握着朝廷存续的物质条件。王敦、苏峻、桓玄与刘裕的时代各不相同，却都说明中枢难以完全脱离外镇与交通节点独立存在。江防为皇统提供了屏障，也不断制造能够要求分享权力的军事中心。

\section{从凉州到龟兹：河西道路上的知识、物资与统治}
\label{sec:eastern-jin-depth-hexicorridor}

河西走廊在十六国时期不是一个只在“西域关系”章节里才出现的边缘。绿洲之间的距离、水源、城郭、驿道和市场决定军队、使节、僧侣与商人能否继续前行。前凉、前秦、后凉、北凉以及北魏先后在这里建立或接收政治权力，但任何一方都不能只靠改元和任命维持通行。灌溉要有人维护，城中粮仓要有人管理，行旅需要供给与保护；这些条件使河西能够连接关中与西域，也使它在中原中心衰弱时拥有重新组织的空间。

吕光从西域返回凉州并据此建立后凉的过程，常被叙述为将领利用前秦崩解的机会。机会本身来自更具体的网络：他所带回的人马、在姑臧能够取得的支持、河西城镇面对前秦权威下降后的选择，以及道路仍可让消息与物资流通。后凉的出现不说明整个走廊都接受同一种统治；不同绿洲、部众与寺院可能与新的军政中心保持不同距离。国号为后来读者提供一个可辨认的标签，不能替代对地方控制程度的逐项判断。

鸠摩罗什的经历让河西道路的文化面向尤其清楚。他从龟兹到凉州、再到长安的移动，既受军事征服和政治安排影响，也依赖不同城市能够提供居处、翻译伙伴与供养。后秦的长安译场不是在空白中突然出现，它接收了西域、河西和关中早已形成的僧侣与文本联系。僧传和经录强调高僧与译本的传承，不能替代沿线所有人的经验；但它们足以说明，战争建立的道路控制有时也会改变知识与宗教资源的流向。

河西的题记、墓葬和文书为这种连续性保留了局部证据。它们可能标出一个人在哪一年任官、一座洞窟由谁供养，或一份契约怎样书写，却不应被放大为整片地区的普查。尤其在政权交替频繁的时期，年号与实际控制范围未必完全重合。材料的分散正提醒我们，地方社会不只是等待中原王朝留下档案；它在自己的城市、寺院与交易中保有可见、也保有难以恢复的时间。

对东晋而言，河西远在北伐军难以持久抵达的地方；对前秦、后秦和北魏而言，它又关系到西向通道、马匹、商贸与象征性的世界联系。不同政权的视角重叠，却没有一个能完全定义走廊的生活。将河西纳入东晋与十六国的叙事，是为了让读者看见，多政权不是只发生在洛阳、长安和建康之间；它同样发生在那些必须靠水源、绿洲和长路维持的区域中心。

\section{家族、官职与地方声望：名门叙事之外的中介}
\label{sec:eastern-jin-depth-local-intermediaries}

东晋史常以王、谢、桓、庾等家族组织人物，确实因为这些家庭在任官、婚姻和军政中留下较多文字。然而门第不是一张脱离地方的名单。一名高门官员的职位要通过僚属、书吏、部曲、佃客、亲属和地方合作者才会变成实际影响；一个家族能够在建康维持声望，也取决于它如何处理籍贯、婚配、葬地、田宅与同乡关系。墓志留下的往往是成功建立声望的表达，不是全部社会关系的透明记录，却能提示我们官职和家族如何被地方化。

迁徙使这种中介作用更为突出。北来者可能以旧有的名望进入建康，也可能在新的地区依赖江南本地人提供土地、婚姻或行政信息。本地精英并非只能被动接受，他们同样能够借朝廷任命、军府联系和市场网络提升位置。两个概念化的集团——“侨姓”与“吴姓”——因此不能替代具体的合作与冲突。不同县、不同世代、不同财力的家庭会有不同策略，传世叙事中关于清谈或礼仪的故事也不应遮住土地、役使与保护的底层条件。

军府是连接都城与地方的另一类中介。将领要有能够记账、传令和筹粮的人，地方官要同军府协调守备与征发，乡里则要在两种权力之间寻求可预测的秩序。当军府强大时，它能够在战争中迅速集结资源；当中枢试图收束兵权时，同一网络又可能成为冲突来源。王敦和桓玄的挑战不只意味着个人野心，也表明他们拥有将人、船和粮从地方调到政治中心的能力。刘裕后来借京口和北府的基础取得优势，亦不能同这种结构分开。

地方声望有时通过宗教和文化活动获得表达。供养寺院、参与译经、刻碑、修墓或主持地方仪式，既可能是信仰与纪念，也可能让家庭在变动中声明自己的关系网。不能把这种行为简化为对政权的忠诚，更不能由少数题记推断全社会的文化取向。它们的价值在于让我们看到：除了诏书、军令和城门，人们还会借文本、仪式和亲属记忆来组织可被承认的身份。

这种中介网络使东晋既有韧性也有局限。它让建康在北方覆亡后仍能找到官员、信息和资源的通路，却使朝廷无法轻易绕过掌握地方关系的人。每次北伐、继承危机或地方叛乱，都是对这些通路的重新测试。读者若只看皇帝与名将，便会把国家能力误作从上而下的命令；若只看地方，又会忽略皇统、官号与军事保护为何仍具有吸引力。两者之间不断协商，才是东晋社会的日常政治。

\section{宗教政策与政治保护：不能用赞助替代信仰}
\label{sec:eastern-jin-depth-religion-politics}

后秦姚兴支持鸠摩罗什译经，北方和江南的王公贵族供养寺院，都是本时期宗教史的重要线索；它们却不能被读成国家对信仰的单向管理。政治保护可以给僧侣、译场和寺院提供安全、经费与城市空间，也会让统治者借文化赞助展示秩序与声望。可是经卷是否被抄写、僧侣如何移动、地方施主为何供养，仍取决于多种关系。宗教活动可以接近宫廷，未必等同于宫廷命令；可以跨越政权边界，也未必因此不受道路与战争限制。

长安译场的材料主要记录高僧、文本和王室关联，读者应避免把它写成城市日常的完整缩影。翻译需要能够理解不同语言的人，也需要纸墨、居处、饮食和抄写；这些物质条件让政治赞助的作用变得可见。另一方面，经卷从长安向河西、江南及其他地区流动时，会经过不同读者、施主与寺院，意义也可能变化。文本的成书、抄写与后世保存各有时间，不能把某一日期固定为所有人同时接受某种教义的时刻。

道教与地方社会的关系亦应作同样处理。孙恩相关叙事中的宗教语言，既是动员与结社可能使用的资源，也是官府标记危险群体的方式。单凭官方称谓，无法决定参与者是因何行动；沿海的经济压力、军役、保护关系和交通条件同样构成背景。把宗教完全从政治中剥离，会忽略它如何进入地方联系；把它完全化为叛乱工具，又会抹去修行、仪式、文本与家庭供养的多样性。

石窟、造像和墓志提供的不是“宗教普及率”，而是不同人愿意留下的公开痕迹。一处供养题记可能让我们知道某人在某地使用何种年号、与哪些人共同投入营造；它并不告诉我们附近没有题名的人是否信仰相同。这样的边界应被保留，因为它使宗教史不再只是朝廷政策的附录，而成为由城市、工匠、施主、道路和文本共同构成的区域实践。

\section{关中与山东之间：远征军到达之后发生什么}
\label{sec:eastern-jin-depth-occupation}

一支军队进入洛阳、广固或长安，史书往往以“克”“入”“降”标记结果。对于治理而言，这些动词只是开始。城门打开以后，统帅要决定旧守军是否留用，地方官能否继续办事，仓粮如何分配，军纪怎样约束，逃散的人何时可能回来。被征服的一方也不是一个整体：有人因旧政权失去保护而寻找新靠山，有人保留武装观望，有人只求让家庭免于再次迁徙。接收若不能进入这些关系，战场上的成功便会很快停在城墙以内。

刘裕围广固与进入关中间隔数年，正可看见不同地区的接收成本。山东靠近淮河和海路，远征军仍须跨越水陆节点并维持围城供给；关中距离江南更远，虽拥有旧都的象征与农业基础，仍需要通过河南道路不断补充兵粮。刘裕的军队能够在一段时间内取得优势，说明东晋军政网络已具有很强的动员能力；长安随后失去，则说明动员能力并不等于在远方建立可自行维持的地方秩序。

后秦的长安本来就不是空城。姚兴时期的官僚、僧侣、市场和军队形成了各自的关系，东晋军到来时必须面对这些遗产。若只以“灭后秦”说明事件，便会忽略新政权要如何处理原有文书、仓储、宗教机构和地方人事。赫连勃勃的夏能够在东晋主力南返后夺取长安，不是因为关中没有居民或行政，而是因为接收尚未把军事控制、地方合作和持续供给牢固地结合起来。不同政权都能利用这座城市的资源，前提是它们能够在通往城市的道路上维持关系。

这种问题在北方各国之间反复出现。前秦吞并前燕后获得河北城镇和人口，但也承继了复杂的军队、地方豪强和运输网络；淝水失利时，这些网络不必再按长安的优先次序行动。北魏向河北扩展亦不是单靠战役结束，而要把平城的军政安排同新附城镇接合。战争可以迅速改变最高权力，接收却是缓慢、常有反复的制度劳动。它的成本由官员、军人、运输者与地方家庭共同承担，而成败也取决于他们是否愿意或能够继续合作。

从这个角度看，东晋北伐的意义不应只按疆域得失计算。每一次北伐都会测试南方能否越过江淮、把粮食和命令送到更远处；也会让北方的地方行动者重新判断建康是否值得依附。祖逖、桓温、谢玄和刘裕处在不同条件中，留下的不是一条单调的成功或失败曲线，而是一组关于道路、驻军和地方接应的经验。后来南北朝的军事组织与统一尝试，正是在这些未被完全解决的经验上继续推进。

\section{流民、安置与再迁徙：迁徙不是一次结束的事件}
\label{sec:eastern-jin-depth-repeated-migration}

永嘉以后的人口移动不能被压缩为一条从北到南的箭头。有人从洛阳、关中或河北出发，抵达江淮后继续向荆州、江南、岭南或海路移动；有人在北方政权之间改换居处；有人因一座城失守而迁入附近坞壁，待局势稍稳又返回耕地。迁徙的节奏取决于亲属、船只、道路、粮食和安全消息，往往包含停留、折返和分离。正史的流亡叙事能够显示大规模危机，不能替每一个家庭画出路线。

安置从来不是只给迁徙者找到住处。朝廷和地方要决定他们归入何种户籍、是否承担赋役、怎样同原有居民划分田地与水源；迁徙者则要依靠同乡、军府、寺院或雇佣关系取得暂时保护。侨州郡保存旧籍与名号，为官员和部分家庭提供制度入口，却不能解决所有人的生产与生活。没有被名册清晰记录的人，可能更依赖军队、地主或地方社群；他们的经历正是材料最难触及的部分。

北方诸政权也面对同样问题。战后俘获、招降、迁徙和屯田都可能改变一地的人口结构，但“迁”在史书中常只是一笔，不足以说明人们是否自愿、抵达何处、能否重新获得土地。前秦扩张把不同区域的人马纳入大军，后赵与燕系政权的战争又让城镇居民反复流离。我们能够确认许多政治转折会带来人口移动，不能把传世的夸张数字变成精确总数。较可靠的写法，是把迁徙同具体的城门、河流、粮道和安置制度相连。

再迁徙改变了地方社会的知识与技能。北来官员带来文书经验，军人带来组织与防卫能力，工匠、商人与僧侣带来不同的生产、交换和文本网络；江南原有的社会也以土地、水利和地方关系吸收或排斥这些资源。把这一过程描述为某一方单向“开发”另一方，会误解双方长期交涉。新的婚姻、任官和市场联系可能由迁徙产生，新的不平等与负担也可能同时产生。

迁徙还解释了为什么国家很难只用行政边界管理社会。一个人在旧籍、实际居所、军府归属和亲属网络之间可能有不同位置；战乱中的家庭更可能在这些位置之间来回移动。东晋朝廷与十六国政权不断试图通过名册、官号和征发使这种流动可见，却无法完全消除它。国家的分类与人们的行动之间的差距，正是多政权时代长期存在的政治空间。

\section{战后城市：市场、寺院与记忆如何回到日常}
\label{sec:eastern-jin-depth-after-war-cities}

城市在战争之后并不会立刻恢复为战前的样子。城墙或许被修补，官署可以重新开门，市场是否有货、居民是否敢回来、道路是否安全，却需要更长时间。建康在王敦、苏峻和桓玄等危机后反复恢复中枢功能，说明它拥有水路、人口与官僚网络的韧性；同样的反复也说明都城周边不断承担驻军、征发和政治斗争的成本。城市能够接住皇统，并不等于城中所有人的生活稳定。

长安、邺城、成都、襄国与姑臧各有不同的恢复条件。长安是关中农业和西向交通的中心，能吸引后秦的译场与后来的军政接收；邺城依托河北城镇与道路，却在后赵崩解中经历剧烈暴力；成都平原给予成汉和东晋接收者资源，但山道限制外部军队；姑臧则将绿洲、市场与走廊交通联结起来。把它们都称为“都城”会遮住这些差异。每座城市重建秩序所需的粮食、劳力、商路与政治保护并不相同。

寺院与墓地是战后城市记忆的特殊场所。寺院可以重新组织供养、抄写和行旅的联系，墓地让家庭在迁徙或安定后声明祖先、婚姻与官职。两类材料都偏向能够投入资源并留下文字的人，不能变成全体居民的声音。它们仍让我们看到，人在政权更替后不只寻求生存，也试图让身份与记忆获得可被承认的形式。城市复苏因而不只是市场重新开张，也是社会关系重新找到可以公开表达的地点。

官方对城市的修复常通过任命、赦令、征粮和重建防务表现出来。这些措施能够说明国家想重新接合什么，不能保证每个地区都以同样速度执行。一个新任太守可能带来文书和军队，也可能触发地方对赋役的新担忧。战后城市的秩序由此总带有临时性：它需要中枢支持，又依赖地方中介；它希望恢复旧的常规，又不得不处理战争留下的人口、债务与权力真空。

读者在国号变化中看到的城市名称，因而应被理解为持续的劳动现场。战争破坏并不使城市的社会功能消失，反而使市场、仓储、寺院、墓地和道路成为下一轮政治重组必须争取的资源。正是在这些地点，东晋与十六国的宏大年序落到可见的生活尺度，也让后来的南北朝继承了不止是疆域，更是许多尚待修补的城市关系。

\section{使节、盟约与消息：多政权间的低强度往来}
\label{sec:eastern-jin-depth-diplomacy}

东晋与北方诸政权并非只在战场相遇。使节、降者、逃亡官员、商旅与僧侣携带消息，在敌对中心之间形成不稳定的往来。正式外交文书会用皇帝、王号和正统安排彼此的位置，实际传递则依赖关隘、渡口、翻译者、护送者和地方能够容纳来人的城市。一次使节到达可证明某种接触发生，不代表双方已经建立稳固同盟；一纸降表也不能保证地方军队和居民会同步转向。

建康尤其需要处理名义与现实的差距。晋室可以以旧朝皇统向北方行动者提出要求，也会利用册封、任命或招抚维持可能的联系。北方政权同样会根据自身军力、都城控制和内部竞争来回应。对于边境的地方集团，接受某一名号有时是获得保护或物资的策略，不必等同于永久的政治认同。正史常把这类选择安排进忠逆框架，读者还应注意其中的道路安全、兵粮和家族生存条件。

消息的速度本身影响决策。淮河、汉水、黄河和河西道路上的一段阻断，可能使都城在数日或更久之后才知道前线改变；当消息抵达时，地方行动者已根据自己的判断结盟或撤离。军事指挥无法摆脱这一延迟，尤其在前秦远征、东晋北伐和燕系政权竞争中更是如此。将领并非在全知条件下选择，而是在有限讯息、季节与资源压力中决定是否推进、驻守或撤回。

这种低强度往来还使文化与行政技术跨越政权边界。官员、书吏和僧侣可以携带文书习惯、经卷、地理知识和人际关系；新的政权会选择吸收、改写或排斥它们。不能把这种流动称为无摩擦的交流：战争、检查、身份怀疑和运输成本随时可能中断联系。它仍说明，多政权不意味着各地完全隔绝。恰恰在边界不断变化的时代，人们更需要依靠能够穿过边界的消息和中介。

\section{灾荒、疫病与史书的缺页}
\label{sec:eastern-jin-depth-disaster}

战争年代的灾荒、疾病与自然灾害常在史书中以简短条目出现：某地饥、某河决、某年有疫。它们不能被视为政治叙事之外的背景。粮食歉收会使军队更难维持，流民增加又会让城市、渡口和地方仓储承受压力；疾病可能影响行军和城内生活，水灾会改变道路与耕作。可是材料通常不足以说明受害范围、死亡数字或每个家庭的应对，不应把简短记载扩写成精确统计。

政权面对灾荒时会赈济、减免、迁民或征调粮食，这些措施首先显示它试图维持何种秩序。诏令的存在不等于救济均匀抵达；地方官是否有仓粮、道路是否可通、军队是否优先取用物资，都会改变结果。东晋和十六国的城市与军镇经常处在战争、迁徙和灾害叠加的状态，单一原因难以解释一次起事或一次政权衰败。保留这种复杂性，比把每场危机归为君主德行或某一群体特性更接近材料能支持的范围。

环境也塑造了不同区域的韧性。江南的水网能够运输粮食，也可能在洪水时破坏耕作；关中依赖农业腹地与灌溉，战乱会使征发与修复更加困难；河西的绿洲对水源与通道极敏感，任何军事争夺都可能影响补给。地方社会会以储粮、迁居、互助、依附或市场交换回应，但我们很少能完整看见这些策略。历史叙事应承认这种缺页，而不是用后世常识假定人们总能以同一种方式渡过危机。

\section{军队中的人：编制、依附与战后去向}
\label{sec:eastern-jin-depth-soldiers}

东晋与十六国的军队并不是只有统帅和兵数。军中有从旧政权转来的士卒、随将领行动的部曲、被征发的地方人、负责运输和修造的人，以及在战后需要重新安置的家属。史书会给出夸张或相互矛盾的数字，难以让我们精确计算其规模；它仍保存了一个事实：军队的形成和解散会持续改变地方的人口、劳动和保护关系。

北府兵的例子尤其重要。它与北来军事人群、京口地区和东晋边防相连，能在淝水等危机中发挥作用；其成员不应被想象成永远同质、只听中枢号令的职业集团。军队需要地方供给、将领的信任和战后奖赏，成功会增强统帅的政治资本，失败则可能使士卒转向其他保护者。刘裕能够从军政网络中上升，便说明军队既是保卫朝廷的工具，也是重组朝廷的力量。

北方政权的军队同样跨越不同地域与身份。前秦的南伐集结来自多处新附地区的人马，规模扩大同时提高协调难度；后赵、燕系和北魏的军事网络也都必须处理部众、城镇守军与农耕人口之间的供给关系。族称或国号不能替代这种组织问题。一支军队能否继续行动，取决于粮食、马匹、道路、亲属安置和上级命令是否还能结合，而非只取决于它在史书中被归入何种人群。

战后去向往往最难恢复。俘虏、降兵、逃散者和伤者可能被编入新军、迁往他地、回归乡里或再次流动；材料通常只记录其中与政治结果有关的一部分。我们不能为他们补造整齐的命运，却应把这种不确定性纳入对战争的理解。每次国号的更替都可能从军队中产生新的地方中介、新的依附关系和新的不安定因素，直到下一场冲突再次将他们卷入。

\section{读法的回望：把事件簇放回同一条时间中}
\label{sec:eastern-jin-depth-reading-return}

本章从南渡、北方政权重组、河西道路、宗教文本、江淮防卫与刘裕北伐分别进入，是为了让不同材料照见同一时期的不同空间。它们不应被读作互不相干的专题。洛阳与长安的陷落使迁徙和侨置成为建康的日常难题；后赵、前燕、前秦与后秦的扩张改变北伐面对的政权组合；河西和长安的宗教网络又让战争控制的道路成为文本和人群流动的条件；淝水的胜负则把远征能力、地方合作与军队组织的差异集中暴露出来。

连续年序的作用，不是把每件事排成不可质疑的因果链。相邻的事件可能共同受迁徙、财政或道路条件影响，却仍需分别说明其具体触发和后果。前秦瓦解没有自动产生北魏，刘裕的军功也没有自动带来南方稳定；每一次转折都经过地方中介、资源分配和新的选择。读者沿着这些事件簇前进，可以同时看见结构性压力怎样累积，也看见行动者在有限信息和物资中仍拥有选择空间。

这也是来源边界改变叙事的地方。正史和编年材料让我们确认称帝、出师、城破和受禅的基本顺序；墓志、题记、经录、城址和文书让我们在局部看见家庭、宗教与地方生活；它们都不足以让我们无缝复原所有人的经验。将可证与不可证并置，并非削弱故事，而是避免用少数著名事件掩盖多政权时代真正的复杂性。东晋与十六国的读者主线，正是在这些可核行动、具体空间与保留的沉默之间向前推进。

\section{技术、手艺与政权的物质能力}
\label{sec:eastern-jin-depth-material-capacity}

战争与行政的能力离不开看似不在帝纪中心的手艺。造船者、铁匠、木工、修城者、驿路维护者和处理文书的书吏，使军队能够渡河、城门能够关闭、征收能够被记录。材料通常把他们压缩在“治城”“造舰”“修器”一类词中，几乎不留下姓名；正因如此，更不应把政权的军事行动写成将领意志自然延伸。每一支远征军的舟楫、马具、粮车和武器，都要通过地方生产、运输和劳役才能出现。

江南水路促进了造船和航运的集中，也使国家在危机中必须争取船户与工匠。刘裕及其他将领能够沿江、沿海调动，不只是因为地图上有河流，而是因为有人熟悉水势、能够修补船只、知道何处可泊何处危险。北方政权则更依赖平原道路、马匹和城镇仓储，但同样需要把铁器、木料与粮食接入军队。不同区域的技术与环境不决定某一政权必然胜利，却规定了它能够怎样行动、在何处付出更高成本。

城市建设也具有政治含义。襄国、邺城、长安、平城和建康的宫城、官署、仓库与防御设施，会把人力和物资集中到新的中心。建筑并不是政权已经稳固的证据，而是它试图使命令、征发和象征持续化的努力。平城的经营把代北军事网络接入可以储粮、设官和接待来者的城郭；建康的维持则依赖江面、街市与周边州县不断提供资源。城的规模与布局可以从考古和后出记述得到局部线索，不能让我们精确复原每次营造或所有劳动者的处境。

技术知识亦随迁徙和道路流动。工匠、书吏、僧侣与商人可能把造纸、书写、建筑、灌溉或航运的经验带到新的地区；政权的接收因而不仅是取得土地和人口，也是在尝试取得让土地与人口发挥作用的知识。河西绿洲的灌溉、巴蜀山道的通行、江南堤防的维护，都不能由一纸中央命令即时复制。地方熟手掌握的信息使他们成为不可忽略的中介，也让新政权必须同旧有社会关系协商。

这些物质条件还帮助理解制度的限度。前秦能在大范围内征调，离不开道路、仓储和城市；当远征过度拉长时，同一系统便暴露协调与补给的脆弱。东晋能以江淮为防线，离不开舟师与水路；当军府控制这些资源时，中枢又难以完全收束兵权。制度得失不只写在诏书和礼仪里，也写在一艘船是否能按时出发、一座仓是否还有粮、一段堤防是否有人修复。将这种层面放入年序，才不会把多政权竞争误作只在宫廷发生的事情。

\section{儿童、老人和照料关系：战争记录之外的家庭时间}
\label{sec:eastern-jin-depth-care}

迁徙和战争不仅改变成年男子的仕宦与军役，也改变家庭中的照料关系。人们离开故乡时，需要带走或安置儿童、老人、病者与不能独立行动的人；一支军队驻进地方，劳力被征发后，留下的人要重新承担耕作、运输和照料。正史很少把这些经验写成连续叙事，墓志偶尔保存的亲属关系也主要来自有资源的家庭。我们不能以想象补齐每一个家庭的故事，却应避免把“南渡”“流民”“迁徙”写成没有年龄、性别和依附关系的人群移动。

家庭分离会影响地方社会如何应对政权变化。有人通过亲属网络寻找过江的船只或新的居处，有人依附军府、地主或寺院以获得暂时保护；婚姻和收养可能成为重建关系的方式，但现存材料不允许把这些可能性泛化为统一模式。对高门家庭而言，谱系、墓志和任官可以维持记忆；对没有文字资源的人而言，照料和依附更可能只以人口流动、劳役与地方传闻的形式留下痕迹。

国家的名册同样看不见完整家庭。户籍、侨置和军籍试图把人归入可征发、可安置的类别，实际居住和亲属关系却可能跨越多个县、军府与政权。战乱中的寡弱者、孤儿与老人最容易落在制度缝隙里；官方救济或赦令若有记载，说明朝廷承认问题存在，却不能保证资源抵达每一处。把这一层不确定性保留在叙事中，有助于读者理解为什么政权更替的代价不会在停战或受禅当天结束。

寺院、家族与乡里有时会成为照料和安置的节点，但不能被浪漫化。宗教机构需要供养，自身也可能受征发和战争影响；家族能够保护部分成员，也可能依赖对外来劳力和佃客的控制；乡里互助同样受资源不足限制。多政权时代的家庭时间因此并不只是私人背景，它与赋役、迁徙、宗教和地方保护交织，构成国家能力最难直接看见、却最持久的基础。

\section{结尾的空间：从建康回望北方，再从北方回望建康}
\label{sec:eastern-jin-depth-spatial-return}

东晋以建康为中心书写自己的存续，北方诸政权则分别以襄国、邺城、长安、龙城、姑臧和平城组织资源。没有一座都城能单独解释这一时期。建康的朝廷需要江淮、荆州、会稽和岭南的船、粮与人；长安的前秦与后秦需要关中腹地、河北、河西和通往西域的道路；平城的北魏需要把代北军事网络接入北方城镇。都城之间并非只有对立，也通过流民、使节、经卷、市场与战争消息彼此牵动。

所以，本章的收束不是把多政权化为一个必然走向统一的前奏。它更关心每个中心如何获得暂时的稳定，又为什么这种稳定会因继承、粮道、地方合作或远征成本而松动。东晋在淝水后保住南方，未能因此接收北方；前秦一度集中广阔资源，未能让所有新附地区在危机中继续服从；北魏从代北发展，靠的是长期把城郭、军队和地方资源重新结合。这些不同路径共同构成南北朝到来之前的历史条件。

对读者而言，最重要的尺度仍是具体空间。渡口决定人能否南下，仓储决定军能否久留，寺院与墓地决定哪些记忆留下，城门决定一夜之间谁获得或失去保护。大事的编年由此不再只是皇帝的顺序，而是这些空间何时被争夺、谁承担代价、何种关系在断裂后被重新接上。东晋与十六国的复杂性，正来自这种不断发生而从不完全完成的接合。
